% AUTO-GENERATED FILE. Source: pub/claims.toml
La evaluación de las explicaciones post hoc en aprendizaje automático sigue
fragmentada entre familias de métricas incompatibles, lo que produce
comparaciones entre métodos sensibles a la elección de los criterios de
evaluación. A partir de un corpus de revisión exploratoria estructurada de 44
artículos, reunido mediante una búsqueda en cinco bases de datos con un
registro de cribado documentado, construimos una taxonomía de cuatro
ejes---objetivo de evaluación, fuente de evidencia, propiedad de calidad y
contexto de la tarea---que organiza el panorama de la medición. A
continuación, realizamos una comparación empírica pareada de LIME y SHAP en 75
celdas emparejadas que abarcan cinco familias de modelos, cinco semillas y tres
tamaños de muestra sobre el conjunto de datos censal UCI Adult, con una
extensión a otros dos conjuntos de datos tabulares. La taxonomía revela tres
brechas recurrentes: las métricas proxy predominan aunque no capturan la
semántica; los constructos basados en el juicio humano siguen insuficientemente
especificados; y la evaluación semántica crece más rápido que su base de
validación empírica. Dentro del protocolo de comparación, los resultados son
consistentes y de gran magnitud: SHAP aventaja a LIME en todos los criterios
orientados a la fidelidad y la estabilidad (estabilidad $d_z{=}3.00$,
fidelidad $d_z{=}4.82$, brecha de fidelidad $d_z{=}2.63$), mientras que LIME
es más parsimonioso y más rápido en las familias de modelos no basadas en
árboles (mediana de 53.3 ms frente a 694.6 ms de SHAP). En los ensambles de
árboles, el orden de latencia depende del modelo: el \texttt{TreeExplainer} de
SHAP fue más rápido que LIME en todas las celdas de XGBoost, pero más lento en
todas las de bosque aleatorio. La extensión a otros conjuntos de datos conserva
el orden de fidelidad---SHAP supera tanto a Anchors como a LIME en todos los
estratos de conjunto de datos y modelo---y muestra que la estabilidad casi nula
de LIME en los datos de Adult no es intrínseca al método: aumenta con el ancho
del kernel y es alta en ambos conjuntos de menor dimensión. En conjunto, la
taxonomía y la comparación empírica motivan un patrón de despliegue
condicionado por la arquitectura del modelo para la clasificación tabular: SHAP
es preferible tanto en fidelidad como en latencia para XGBoost, LIME conserva
una ventaja de latencia en arquitecturas de caja negra sin un explicador
específico del modelo, y en otros ensambles de árboles la latencia debe medirse
antes de elegir. Todos los artefactos experimentales están disponibles
públicamente.
